\documentclass[conference]{IEEEtran}
\IEEEoverridecommandlockouts

\usepackage{graphicx}
\usepackage{booktabs}
\usepackage{amsmath}
\usepackage{cite}
\usepackage{url}

\title{Reproducible Agent-Assisted Detection and Forensics in a Student Cloud Security Range}

\author{\IEEEauthorblockN{Student A, Student B, Student C, Student D}
\IEEEauthorblockA{Centre for High Performance Computing / SCC26 Student Project\\
South Africa\\
\{student1, student2, student3, student4\}@example.invalid}
\and
\IEEEauthorblockN{Mentor Name}
\IEEEauthorblockA{Mentor / Affiliation\\
mentor@example.invalid}}

\begin{document}
\maketitle

\begin{abstract}
% 100--150 words. State: problem/context; what you built/tested; method; one or two
% evidence-backed findings; why the result matters. Do not put literature review here.
This paper presents ...
\end{abstract}

\begin{IEEEkeywords}
OpenStack, Kubernetes, GitOps, reproducibility, % add 2--4 project-specific keywords
\end{IEEEkeywords}

\section{Introduction}
% ~0.5 page.
% 1) What problem are you studying?
% 2) Why does it matter?
% 3) What gap/question does this student project address?
% 4) End with 2--4 concrete contributions, not promises.

\section{System Architecture and Context}
% ~0.5--0.7 page. Describe only the components needed to understand your experiment.
% Include one readable architecture figure. Separate Sebowa OpenStack/Kubernetes from
% externally allocated A100/H200 resources when applicable.

\begin{figure}[t]
  \centering
  % \includegraphics[width=\columnwidth]{figures/architecture.pdf}
  \fbox{\parbox[c][30mm][c]{0.92\columnwidth}{\centering Replace with architecture figure}}
  \caption{System architecture and experimental boundary.}
  \label{fig:architecture}
\end{figure}

\section{Methodology}
% ~0.7--1.0 page.
% Define reproducible procedure, versions, resources, workload/scenario, variables,
% controls, metrics and success/failure criteria. State what was held constant.

\subsection{Platform and Software Baseline}
% Git SHA(s), immutable image digest(s), Kubernetes/OpenStack resources, relevant versions.

\subsection{Experimental Procedure}
% Exact experiment sequence and number of repeats where meaningful.

\subsection{Metrics and Evidence}
% Name metrics/log sources and how timestamps/results were correlated.

\section{Results}
% ~1--1.3 pages. Present measured observations before interpretation.
% Use compact tables/plots. Do not report fake placeholders in the final paper.

\begin{table}[t]
\centering
\caption{Example result structure (replace with real measurements).}
\label{tab:results}
\begin{tabular}{lll}
\toprule
Configuration & Metric & Measured value \\
\midrule
A & metric-1 & TBD \\
B & metric-1 & TBD \\
\bottomrule
\end{tabular}
\end{table}

\section{Discussion and Limitations}
% Explain why results behaved as observed; relate to architecture/resource constraints;
% discuss measurement uncertainty, failed tests, threats to validity and what cannot be
% concluded from this small POC.

\section{Reproducibility}
% Short but explicit: repository revision, image digests, protected inputs, rebuild test,
% what was automated, what remained manual, and whether the key result was reproduced.

\section{Conclusion and Future Work}
% 1 short paragraph: answer the research question at the level supported by evidence.
% Then name the most defensible next step(s), not a wishlist.

\section*{Acknowledgment}
% Acknowledge CHPC/NICIS/CSIR, mentors, infrastructure providers and collaborators as appropriate.

% Project-specific minimum evidence ideas retained as author comments:
% - At least one approved scenario producing both host and network evidence.
% - Wazuh and Suricata detections correlated on a common timeline.
% - A human-reviewed ACP/Hermes incident explanation grounded in retained evidence.
% - Documented false positives, false negatives and telemetry gaps.
% - A clean reset/replay path that does not require attacking any system outside the authorised ranges.

\bibliographystyle{IEEEtran}
\bibliography{references}

\end{document}
